\section{Query-conditioned evidence fusion}
\label{sec:method}
The method forms a response in three stages: determine which observations
apply to the query, assign their spatial support and exclusion authority, and
combine their admitted responses through geometry. Semantic and complete-instance
outputs retain their own objectives and readouts. The following equations
describe the implemented system and identify its query-specific realizations.

\begin{figure}[ht]
\centering\includegraphics[width=\linewidth]{figures/method_overview.pdf}
\caption{Our evidence-to-response method. Applicability and spatial authority
determine how native evidence acts on the response. Geometric aggregation and
complete membership have distinct typed implementations. Actual selections are
evaluated separately from continuous semantic responses.}
\label{fig:method}
\end{figure}

\subsection{Native evidence and query applicability}
For query $q$, observations $\mathcal E_q$ supply responses $y_e\in[0,1]$,
nonnegative geometric operators $A_e$ and visible domains $v_e$. Nonempty
operator rows are normalized. Text similarities and local visual correspondence
stay in their native teacher spaces; points and reference masks retain their own
semantics. We combine response evidence, not raw feature vectors. LERF, NVOS
and ScanNet use declared source-assisted observations; UQIS reads its map without
source RGB or depth at query time.

Let $g_e(q)$ express whether observation $e$ applies to the requested target.
For text candidates, we compare one- and two-component native-score models
within candidates already detected for the same query. If their Bayesian
information criteria are $B_1,B_2$, the implemented support is
\begin{equation}
 a_e=(1-\pi_2)+\pi_2\gamma_{\rm high}(s_e),\quad
 \pi_2=\sigma((B_1-B_2)/2),\quad g_e=\mathbf1[a_e\geq0.5],
 \label{eq:admission}
\end{equation}
where $\gamma_{\rm high}$ is responsibility for the higher-score component.
A single mode supplies no relative identity veto. Degenerate or missing evidence
does not invent one. Rejected regions become unknown rather than background.
Mixture support is not a calibrated probability of semantic truth.

This rule has a controlled LERF result. Reference queries instead inherit
identity through a positive-prompt mask and source tracks; categories combine
native probabilities with literal-category evidence. Image queries use local
correspondence to select identity hypotheses. Actual points supply explicit
spatial conditions. These are typed estimators, not unreported copies of
Equation~\ref{eq:admission}.

\subsection{Spatial authority and geometric response formation}
An applicable observation need not describe complete extent. Let $p_e,n_e$
denote disjoint support and exclusion domains and $c_e^+,c_e^-$ their weights:
\begin{equation}
 w_e=v_e g_e(q)(c_e^+p_e+c_e^-n_e).
 \label{eq:authority}
\end{equation}
Unobserved space and rejected detections carry no exclusion solely because they
lack a match. Predicted reference-mask complements and category competition can supply
exclusion under their stated completeness assumptions. These are fallible
provider estimates, not hard negative user prompts. Our LERF path vetoes negatives where
native positive support exists. NVOS uses interior positive/negative regions and
tracked negative reliability. ScanNet keeps both native priors and source weights.

For direct consensus, define $m=\sum_e A_e^\top w_e$ and
$b=\sum_e A_e^\top(w_ey_e)$. The response is
\begin{equation}
 r_i=\begin{cases}
 (b_i+\lambda_i\bar r_i)/(m_i+\lambda_i),&m_i>10^{-8},\\
 r_{0,i},&\text{otherwise}.
 \end{cases}
 \label{eq:consensus}
\end{equation}
LERF uses $\lambda=0$. ScanNet uses a two-native mean $\bar r$ and prior strength
0.1 times visible mass. This is an established weighted-consensus operation;
the tested design concerns admitted evidence and its actual use in the response.
The active ScanNet implementation has no fitted local residual head.

NVOS instead minimizes robust rendered discrepancy,
\begin{equation}
 \min_{r\in[0,1]^N}\sum_e\langle w_e,\rho_{0.25}(A_er-y_e)\rangle,
 \label{eq:reference}
\end{equation}
with 120 mass-preconditioned updates from a neutral 0.5 response. Its output
adapter zeros unsupported elements. Changing this solver while preserving
evidence is an ablation. Our controls do not support replacing every typed
operator by one inverse solver.

\subsection{Complete extent and typed outputs}
For source-assisted actual selections, our implementation also fits membership
against the silhouette of the selected Gaussians. This is distinct from rendering
their semantic scores in the full scene. Let $z_i\in[0,1]$ be membership and
$o_{vi}(u)$ the uncomposited per-Gaussian opacity at pixel $u$. The selected
silhouette has the compositing form
\begin{equation}
 \alpha_v(u;z)=1-\prod_i(1-z_i o_{vi}(u)).
\end{equation}
The implemented loss combines foreground/background-balanced binary
cross-entropy and soft Dice on admitted source pixels, plus a 0.1-weighted
projected-footprint membership penalty. The footprint is normalized by its
cost at source-initialized membership. Optimization uses at least 240 Adam
updates (extended to the source-view count when larger) at learning rate 0.05,
cycling through the active views. Actual membership is then
thresholded at 0.5 and its selected opacity is rendered. The existing projection
readout thresholds quantized 8-bit alpha above 10. This describes an existing
v10 capability, not a new solver introduced by the subsequent NVOS experiment.

\paragraph{Native image identity and conditional complete membership.}
For an image query, let $q_u,f_v$ be normalized query and regional patch
features, with unit-sum weights $a_u,b_v$. Local correspondence uses the
symmetric weighted Chamfer score
\begin{equation}
 k(Q,R)=\tfrac12\left[\sum_{u\in Q}a_u\max_{v\in R}q_u^\top f_v+
                         \sum_{v\in R}b_v\max_{u\in Q}q_u^\top f_v\right].
\end{equation}
The deployed image route averages the candidate ranks from native SigLIP
matching and local DINOv2 correspondence. It selects a nonempty region, then
uses its actual member nearest the coordinate-wise 3D median as the identity
anchor. Both teacher spaces remain unchanged. An actual point supplies a
geometric anchor directly; an image-derived anchor remains an identity estimate.

For anchor $a$ and every carrier element $i$, the complete-membership decoder
predicts
\begin{equation}
 p_i=f_\theta(\phi(i,a)),\qquad S_a=\{i:p_i>\tau\},\quad p_a=1.
\end{equation}
The 15-dimensional evidence vector retains view-balanced nested-mask votes,
any-mask/all-mask agreement and their disagreement, common visibility and its
presence flag, anchor/element visibility fractions, bounded signed displacement
and distances, native appearance agreement, and native-support availability.
Nested masks share a view's vote instead of becoming independent observations.
The decision is evaluated over the complete carrier, without requiring an
intermediate mask to be a pure or complete instance.

The deployed decoder is an Extra Trees classifier with 64 trees, at most 512
leaves per tree and minimum leaf size 32. It is inherited from training on
independent 3D instance labels in 22 real-observation rooms; three separate
calibration rooms select $\tau=0.35$. Neither split overlaps the benchmark
rooms. This supervised decoder is disclosed as part of the method; freezing
teacher features does not make the complete system training-free. The same
anchor-positive rule is applied to inferred image anchors. Its identity
condition is fallible, and an ablation removing only that assignment leaves
all measured selected sets unchanged.

Continuous semantic response, category argmax and actual instance selection
have different readouts. LERF and NVOS maintain independent Gaussian sets;
changing the semantic response does not change that set automatically.
Selected-set projection IoU uses 2D GT, not volumetric GT. UQIS complete
membership is not Equation~\ref{eq:consensus}. These distinctions are essential
to evaluating functional improvement.

\paragraph{Why the optimization target and readout must agree.}
Removing unselected Gaussians changes ray transmittance, not just the values
averaged by the original renderer. A simple ray shows the distinction: an
unselected front Gaussian with opacity 0.9 followed by a selected Gaussian with
opacity 0.2 gives a full-scene normalized member response
$0.02/(0.9+0.02)\simeq0.0217$. Rendering the selected set alone instead gives
alpha 0.2. The two outputs disagree under their existing thresholds despite
identical membership. This is an illustrative counterexample to treating member
scores as calibrated semantic responses, not an accuracy result. Accordingly,
functional inheritance preserves evidence, optimization target and intended
readout together. Our field-by-readout experiment tests this distinction on
retained NVOS fields.

\paragraph{Image-conditioned response refinement.}
For an observed view $t$, let $z_t^0=A_t s^0$ be the rendered v10 response.
We use its thresholded extent as a box and its bounded log-odds as a dense
prompt for a frozen image decoder $D$:
\begin{equation}
 p_t=D(I_t,\operatorname{box}[z_t^0\geq0.5],
       \operatorname{resize}[\operatorname{logit}(\operatorname{clip}(z_t^0))]).
\end{equation}
Clipping uses $[10^{-4},1-10^{-4}]$; the dense prompt is resized to the
decoder's actual $288\times288$ grid. The box and dense mask are fallible
predictions, not extra annotations. Single-mask decoding avoids target-label
candidate selection. Paired reconstructed/observed RGB controls isolate image
access under this fixed interface. Sparse-point and unboxed controls test
whether the inherited response preserves identity while restricting extent.

The coarse control uses confidence-weighted area pooling of $p_t$ on a
$512\times384$ grid. Native-ray writeback instead retains $y_t=p_t$ on the
$1008\times756$ output grid and uses $w_t=|2p_t-1|v_t$. Here $v_t$ admits rays
with nonzero observed mass and more than half of their rendered alpha supplied
by the original observed Gaussians. The normalized observation matrix
$\bar A_t$ is formed on the same ray grid, using the fixed original Gaussian
geometry. Initialize at $s^0$ and take 120 bounded Huber updates for
\begin{equation}
 \min_{0\leq s\leq1}\sum_u w_{tu}\rho_{0.25}((\bar A_t s)_u-y_{tu}),
 \qquad s_i=s_i^0\ \text{when no admitted observation supports }i.
\end{equation}
We measure both $p_t$ and the actual re-rendered 3D response. Geometry and
teacher weights remain fixed. Because this update uses view $t$'s image,
the result is view-conditioned; a query-global field and preserved source-view
accuracy do not follow from target-view gains. Original source-teacher
residuals are reported as a separate compatibility check.

\paragraph{Completing query scope before transport.}
The current NVOS extension builds a complete reference response before cross-view
transport. For reference image $I_r$, let $\mathcal P,\mathcal N$ be original
positive and negative scribble patches. A patch containing both signs is unknown;
unmarked patches do not supply negative labels. With frozen normalized DINOv2
descriptors $f_r$, define
\[
 h_r(u)=\max_{v\in\mathcal P} f_r(u)^\top f_r(v)
       -\max_{v\in\mathcal N} f_r(u)^\top f_r(v),
 \qquad a_r(u)=\sigma(h_r(u)/\tau_r).
\]
The source-only scale $\tau_r$ is the median signed leave-one-out reference
margin divided by $\log9$. Missing signed support or a nonpositive scale falls
back to the original reference mask. Upsample $a_r$ and use the same dense-mask
and extent-box interface to obtain a semantic reference response from frozen
SAM3. A second candidate additionally uses 16 deterministically spread points
per original scribble sign. Threshold both at 0.5; for the constrained candidate,
impose all original positive/negative pixels exactly, without relabelling
unknown pixels. Agreement with sparse annotations does not prove complete scope.

Two fixed reference-only admission policies are compared. The feasibility policy
retains the semantic candidate when it satisfies every original signed pixel,
otherwise choosing the constrained candidate. The source-quality policy chooses
\[
 \hat k=\arg\max_{k\in\{\mathrm{semantic},\mathrm{constrained}\}}
             \widehat Q_r(k),
\]
where $\widehat Q_r$ is SAM3's predicted mask quality at reference decoding,
before the final hard annotation constraints. Ties prefer the semantic candidate;
missing quality falls back to feasibility. These scores are fallible predictions,
not target GT or calibrated correctness probabilities. Neither rule reads a
scene identifier, target score or task-specific threshold.

The selected whole reference mask initializes frozen two-frame SAM3 memory,
producing $y_t=T(I_r,m_r^{\hat k},I_t)$. The same native-ray fit above then
creates and actually re-renders the 3D response. Original negative scribbles,
DINOv2 and observed target RGB remain declared extensions to v10's input contract.
This changes response construction and admission, keeping teacher weights and
Gaussian geometry fixed. Both fixed policies are reported on all eight tasks;
reference-admission design follows the opened-cohort completion diagnosis.
Earlier geometric region-admission and paired-residual alternatives remain
matched diagnostics, rather than additional required stages of this pipeline.

\paragraph{Optional complete-response replacement.}
A subsequent extension generates intact alternatives and predicts relative task
utility from query/response signatures. A calibrated margin retains the
incumbent or returns one complete alternative. Its L11 variant uses a structural
potential plus a query-conditioned residual. This extension is evaluated
separately: its small image gain and unfinished LERF/NVOS transfer do not account
for the historical component gains.
